% Appendix. Source files for every number here (2026-07-31):
%   labeling params    meta field of /mnt/scratch/lh/labels/final3_env_rc_*.npz
%                      and final2_ol_*.npz
%   sigma sweep        final3_env_sig{025,100}_rc_*.npz, recomputed 2026-07-31
%   policy action RMSE results/rmse/*.json (diffusion policies)
%                      results/robocasa/rmse_groot.log (GR00T N1.5)
%   predictor          results/predictor/rc_seq_v3/S_<task>_h<1|4>_s<0|1|2>/
%                        {metrics,calib}.json
%   online thresholds  results/predictor/rc_seq_v3/online_thresholds.json
%   labeling wall time DONE lines of results/robocasa/label3_*.log
%   amplification /       recomputed from final3_env_rc_*.npz on 2026-07-31:
%   lambda_full agreement   delta reproduces Table 1 to 4 dp; e^{16 lambda} at
%                           delta / p99 / max; spearman(lambda_task, lambda_full)

\appendix

\section{Error Growth Over the Execution Horizon}
\label{app:theory}

This appendix expands the premise paragraph of Section~\ref{sec:method}: the
statements we take from the literature, the two lines of algebra that specialize
them to one executed chunk, and what our measurement of $\lambda$ does and does
not license.

\paragraph{What is quoted.}
Three results are used, none of them ours.

\emph{(i) Compounding from distribution shift.} \citet{ross2010efficient},
Theorem 2.1:
if $\hat\pi$ has expected $0/1$ loss at most $\varepsilon$ under the expert's
own state distribution, then $J(\hat\pi) \le J(\pi^\star) + T^2\varepsilon$ over
a horizon of $T$ steps, and the quadratic is tight, since a system, cost and
policy can be constructed that attain it. This is the distribution-shift form of
compounding error and it assumes nothing about the dynamics.

\emph{(ii) The continuous-action obstruction.}
\citet{simchowitz2025pitfalls} define $f$ to be $(C,\rho)$-exponentially
incrementally input-to-state stable when \eqref{eq:eiss} holds for all pairs of
initial states and input sequences, with $C\ge 1$ and $\rho \in [0,1)$. Their
Lemma 2.1 gives the benign direction: if the \emph{closed loop} $(f,\hat\pi)$ is
$(C,\rho)$-stable and $\hat\pi$ is $L_{\hat\pi}$-Lipschitz, the execution error
is at most $\frac{C}{1-\rho}(2+L_{\hat\pi})$ times the expert-distribution
error, with no factor of the horizon at all. Their main theorem gives the
obstruction: there are imitation problems with smooth deterministic experts and
stable dynamics on which any smooth, Markovian, simply-stochastic imitator
suffers execution error $\exp(\Omega(T))$ times its training error. Closing the
loop is not by itself enough; the closed loop has to be stable, and learning
from expert data alone does not deliver that.

\emph{(iii) Chunking as the remedy.} \citet{zhang2025chunking} analyze
action chunking directly. Their Proposition 3.1 states that if the learned pair
$(\hat\pi,\hat f)$ is $(C_{\text{ISS}},\rho)$-stable and the chunk length
$\ell$ exceeds $\log(1/\rho)^{-1}\log(\mathrm{poly}(L,C_{\text{ISS}}))$, then the
chunked policy on the true dynamics is again exponentially stable, with constant
$\log(1/\rho)^{-1}\mathrm{poly}(L,C_{\text{ISS}})$ and rate $\rho^{1/2}$; their
Theorem 1 then bounds trajectory error by a constant multiple of regression
error, with the constant free of $T$. The mechanism is that a long enough
open-loop chunk lets the plant's own contraction wash out the mismatch between
the learned and true dynamics before the next replan is committed. Note the
direction: this is a lower bound on $\ell$, and it presumes $\rho<1$.

\paragraph{What is derived.}
The bounds in \eqref{eq:chunk} and~\eqref{eq:window} follow from
\eqref{eq:eiss} by substitution and are not theorems from those works.

For the labeling probe, both branches start at the same state, so the first term
of \eqref{eq:eiss} vanishes, and only $u_1$ differs, so the sum has one nonzero
term:
$\lVert x_j - x'_j\rVert \le C\rho^{\,j-1}\lVert u_1 - u'_1 \rVert
= C\sigma_u\,e^{(j-1)\lambda}$ with $\lambda := \log\rho$. Taking logarithms
gives an affine function of $j$ with slope $\lambda$, which is what
Algorithm~\ref{alg:label} fits.

For one executed chunk, the action error is present at every step because no
replan intervenes, so $\lVert u_i - u'_i\rVert = \varepsilon$ for all $i \le k$
and the sum is geometric,
$C\varepsilon\sum_{i=0}^{k-1} e^{i\lambda}
= C\varepsilon\,(e^{\lambda k}-1)/(e^{\lambda}-1)$,
which is \eqref{eq:chunk}. Its three branches are: for $\lambda<0$ the sum is
increasing in $k$ and converges, so $C\varepsilon/(1-e^{\lambda})$ bounds it for
every $k$; at $\lambda=0$ the ratio has a removable singularity with value $k$,
giving $C\varepsilon k$; and for $\lambda>0$ the numerator dominates, leaving
$C\varepsilon e^{\lambda k}/(e^{\lambda}-1)$.

For the episode, write $e_n$ for the deviation at the $n$th chunk boundary and
model one replan as a contraction by $\gamma$, so that
$e_{n+1} \le \gamma\!\left(e^{\lambda k} e_n
+ C\varepsilon\,\frac{e^{\lambda k}-1}{e^{\lambda}-1}\right)$.
This is a linear recursion with gain $\gamma e^{\lambda k}$, bounded over
$n = T/k$ chunks exactly when $\gamma e^{\lambda k} < 1$, that is when
$k < \log(1/\gamma)/\lambda$. Combining it with the lower bound of (iii) gives
the admissible window \eqref{eq:window}.

\paragraph{The threshold in amplification units.}
The automatic threshold $\delta$ of Section~\ref{sec:method} is set from the
$\lambda$ distribution alone and never sees a chunk length, so where it lands in
the units of \eqref{eq:chunk} is a check rather than a construction.
Table~\ref{tab:amp} converts it. On all four RoboCasa tasks the threshold sits
where one 16-step chunk multiplies an action error by between 2.5 and 3.9, the
top percentile of stamps sits at 4.5 to 6.9, and the worst stamps reach 19 to
107. Equivalently, holding the amplification at a stamp on the threshold to a
factor of two would require chunks of 8 to 12 steps rather than 16.

\begin{table}[h]
\centering
\small
\begin{tabular}{lcccccc}
\toprule
& \multicolumn{2}{c}{threshold} & \multicolumn{2}{c}{99th pct.} & \multicolumn{2}{c}{worst stamp} \\
\cmidrule(lr){2-3} \cmidrule(lr){4-5} \cmidrule(lr){6-7}
Task & $\delta$ & $e^{16\delta}$ & $\lambda$ & $e^{16\lambda}$ & $\lambda$ & $e^{16\lambda}$ \\
\midrule
OpenCabinet & .085 & 3.9 & .115 & 6.3 & .222 & 35 \\
OpenDrawer & .076 & 3.4 & .121 & 6.9 & .223 & 36 \\
PickPlaceToCabinet & .065 & 2.8 & .095 & 4.5 & .184 & 19 \\
TurnOnSinkFaucet & .058 & 2.5 & .098 & 4.8 & .292 & 107 \\
\bottomrule
\end{tabular}
\caption{The measured $\lambda$ distribution read as the per-chunk
amplification factor $e^{\lambda k}$ of \eqref{eq:chunk} at the executed chunk
length $k{=}16$. The threshold column is the automatic $\delta$ of
Table~\ref{tab:labels}, which is computed from the below-median half of the
$\lambda$ distribution and has no knowledge of $k$.}
\label{tab:amp}
\end{table}

\paragraph{What the correspondence does not license.}
Four caveats, in decreasing order of how much they matter.

\emph{$\lambda$ is measured on an output, not on the state.} The left-hand side
of \eqref{eq:eiss} is $\lVert x_j - x'_j\rVert$, the full state. Our label is
the growth rate of a task metric $y = h(x)$, and composing a Lipschitz $h$ with
\eqref{eq:eiss} bounds $\lVert y_j - y'_j\rVert$ by $L_h C \rho^{j-1}\sigma_u$,
so the fitted slope estimates the rate along the observed coordinates rather
than the state's leading exponent. The two are not interchangeable here. The
labeler also records the slope on the full simulator state, and on RoboCasa the
two are nearly uncorrelated: Spearman $0.21$, $0.07$, $-0.02$ and $0.01$ on the
four tasks, with the full-state slope having a \emph{negative} median
($-0.019$ to $-0.010$) where the task-metric slope is near zero or positive.
The full state settles while the coordinates the task cares about separate. We
use the task metric deliberately, because whether the chunk still lands where
it was aimed is a statement about those coordinates, but the label should be
read as an output-space expansion rate and not as the plant's leading Lyapunov
exponent.

\emph{$\gamma$ is assumed, not measured.} Nothing in this paper measures the
contraction a single replan buys. The closed-loop labels of
Section~\ref{sec:method} measure whether replanning helps at all at a state,
which is the sign of the effect, not its magnitude, so \eqref{eq:window} is a
statement of shape rather than a formula we can evaluate. Measuring $\gamma$ is
the natural way to turn the switching rule from two levels into a schedule.

\emph{$C$ and $L$ are not measured either.} The lower bound in
\eqref{eq:window} therefore cannot be evaluated for our systems, and we do not
claim $k{=}16$ satisfies it. What the window says is that the two bounds move
independently: the lower one is a property of the system and the learned model,
the upper one of the state.

\emph{The rate is finite-time.} $\lambda$ is fitted over a $K{=}24$ step window
and is a finite-time exponent in the sense of \citet{haller2015lcs}, not an
asymptotic one. Figure~\ref{fig:labeling}(b) shows that neither divergence curve
is a clean straight line; the slope summarizes the window, and
\eqref{eq:chunk} inherits that approximation.

\section{The Labeling Procedure in Full}
\label{app:labeling}

Algorithm~\ref{alg:label} is the procedure of Section~\ref{sec:method} written
out. It is the whole method: there is no training step and no policy query.

\begin{algorithm}[h]
\caption{Open-loop stability label at one state}
\label{alg:label}
\begin{algorithmic}[1]
\Require simulator state $s_{t_0}$, recorded actions $a_{t_0},\dots,a_{t_0+K-1}$,
         horizon $K$, branches $N$, scale $\sigma_u$
\State restore the simulator to $s_{t_0}$
\For{$j = 0$ to $K-1$}\Comment{nominal branch}
  \State step the simulator with $a_{t_0+j}$; record $y_j \gets$ task metric
\EndFor
\For{$n = 1$ to $N$}
  \State $\tilde a \gets a_{t_0}$; \; $\tilde a[\text{translation}] \mathrel{+}= \sigma_u \varepsilon$,
         \; $\varepsilon \sim \mathcal{N}(0, I_3)$
  \State restore the simulator to $s_{t_0}$
  \For{$j = 0$ to $K-1$}\Comment{perturbed branch: same actions after the first}
    \State step with $\tilde a$ if $j = 0$ else $a_{t_0+j}$; record $y^{(n)}_j$
    \State $d^{(n)}_j \gets \lVert y^{(n)}_j - y_j \rVert$
  \EndFor
\EndFor
\State \Return $\lambda \gets$ OLS slope of
       $\frac{1}{N}\sum_n \log \max(d^{(n)}_j, 10^{-12})$ against $j$
\end{algorithmic}
\end{algorithm}

\paragraph{Settings.}
Every dataset in this paper uses $K{=}24$, $N{=}8$, a stride of 2 between
labeled stamps, and perturbs the three translation dimensions only. The gripper
dimension is never perturbed, because a perturbation there changes which
discrete contact mode the branch enters and measures a different thing.

\paragraph{Task metric.}
The metric $y$ is the quantity whose divergence defines the label, and it
differs by platform because the observation spaces do. On robomimic, MimicGen
and LIBERO it is the end-effector position concatenated with the task object's
observation vector. On RoboCasa it is the end-effector position concatenated
with the gripper joint positions, because the kitchen environments have no
single distinguished object observation. Both are norms over quantities in
metres, so $\lambda$ carries units of inverse control steps in both cases.
The change of metric between platforms means $\delta$, and therefore the
unstable fraction, is not directly comparable across the two groups; the
within-episode structure that Claim 1 rests on is.

\paragraph{Action ordering.}
The RoboCasa demonstrations are stored in a different action order from the one
the controller consumes:
\begin{center}
\small
stored $[\text{base}_3, \text{torso}_1, \text{mode}_1, \text{arm}_6, \text{grip}_1]$
\quad$\rightarrow$\quad
consumed $[\text{arm}_6, \text{grip}_1, \text{base}_3, \text{torso}_1, \text{mode}_1]$.
\end{center}
An earlier version of our labels replayed the stored order directly, which
drives the arm with base commands and makes the nominal branch itself a failure.
Replay validation caught it: the remapped order reproduces the demonstration
end state to within 4 to 7 centimetres, the raw order misses by about a metre.
Every RoboCasa number in this paper comes from labels regenerated after the fix.
We report this because a silent action-order mismatch produces labels that look
plausible, with a sensible $\lambda$ distribution and a sensible threshold, and
is invisible without an explicit replay check.

\paragraph{Determinism.}
The label is reproducible to the last digit only if the simulator call sequence
is reproduced exactly, including calls that read nothing. Re-running the
procedure while omitting a redundant \texttt{reset\_to} shifts MuJoCo's
warm-start cache, which perturbs the constraint solver and moves $\lambda$ by a
few thousandths, enough to flip stamps near $\delta$. Regenerating the
divergence curves of Figure~\ref{fig:labeling} with the call sequence
reproduced exactly recovers the stored labels with a maximum absolute
difference of $0$ over the 25 stamps checked.

\section{The Perturbation Scale $\sigma_u$}
\label{app:sigma}

\paragraph{How it is set.}
The intent stated in Section~\ref{sec:method} is that $\sigma_u$ should be the
trained policy's validation action RMSE, so that the injected perturbation is
the size of the error the policy actually makes. Table~\ref{tab:sigma} gives
the value used for every dataset. For the diffusion-policy platforms this rule
is applied directly, using the position-dimension RMSE of the trained
checkpoint on held-out demonstrations. For RoboCasa it is not: all four tasks
use $\sigma_u{=}0.05$, chosen before the GR00T checkpoint was evaluated and
kept for consistency across the sweep.

\begin{table}[h]
\centering
\small
\begin{tabular}{llc}
\toprule
Platform / dataset & $\sigma_u$ & source \\
\midrule
RoboCasa, all four tasks & 0.05 & fixed, see below \\
robomimic can & 0.225 & policy RMSE \\
robomimic lift & 0.165 & policy RMSE \\
robomimic square & 0.220 & policy RMSE \\
robomimic tool\_hang & 0.164 & policy RMSE \\
MimicGen coffee preparation & 0.114 & policy RMSE \\
MimicGen nut assembly & 0.103 & policy RMSE \\
MimicGen square & 0.220 & policy RMSE, square ph \\
MimicGen stack & 0.203 & robomimic mean \\
LIBERO-10, all ten tasks & 0.203 & robomimic mean \\
\bottomrule
\end{tabular}
\caption{Perturbation scale per dataset. Policy RMSE is the position-dimension
action RMSE of the trained diffusion policy on held-out demonstrations. Two
datasets use the mean over the robomimic policies because their own policy was
not yet trained when the labels were produced.}
\label{tab:sigma}
\end{table}

\paragraph{The RoboCasa gap.}
Applying the stated rule to RoboCasa would give a much larger perturbation than
the one we use. The GR00T N1.5 checkpoint's position-dimension action RMSE on
held-out demonstrations is 0.565 on OpenCabinet, 0.426 on OpenDrawer, 0.479 on
PickPlaceCounterToCabinet and 0.299 on TurnOnSinkFaucet, pooling to 0.467, an
order of magnitude above the 0.05 used for labeling. The sensitivity analysis
below covers a fourfold range around 0.05 and does not reach the rule-implied
value, so it does not by itself justify the choice at that scale.
\pending{relabel at least one RoboCasa task at the measured RMSE and report
whether the regime structure survives; if it does not, the honest statement is
that $\sigma_u$ is a free parameter of the label rather than a policy-derived
one}

\paragraph{Sensitivity.}
Table~\ref{tab:sigmasweep} relabels all four RoboCasa tasks at
$\sigma_u \in \{0.025, 0.05, 0.10\}$. The threshold $\delta$ moves with the
probe, as it must, since a larger kick produces larger separations everywhere.
What matters is that the per-stamp ordering of $\lambda$ is nearly preserved
and the resulting unstable flag barely moves: adjacent settings agree on 96.6
to 99.4 percent of stamps. The structure the paper relies on is therefore a
property of the dynamics, not of the probe size, over this range.

\begin{table}[h]
\centering
\small
\setlength{\tabcolsep}{4pt}
\begin{tabular}{lccccccc}
\toprule
& \multicolumn{3}{c}{$\delta$ at $\sigma_u$} & \multicolumn{3}{c}{Spearman $\rho$ of $\lambda$} & agree \\
\cmidrule(lr){2-4} \cmidrule(lr){5-7} \cmidrule(lr){8-8}
Task & .025 & .05 & .10 & \tiny .025/.05 & \tiny .05/.10 & \tiny .025/.10 & \% \\
\midrule
OpenCabinet & .100 & .085 & .067 & .98 & .96 & .90 & 99.1--99.4 \\
OpenDrawer & .088 & .076 & .061 & .96 & .93 & .85 & 99.0--99.3 \\
PickPlaceToCabinet & .078 & .065 & .046 & .97 & .95 & .89 & 96.6--98.4 \\
TurnOnSinkFaucet & .068 & .058 & .046 & .95 & .93 & .82 & 97.2--98.3 \\
\bottomrule
\end{tabular}
\caption{Relabeling all four RoboCasa tasks across a fourfold range of
$\sigma_u$, on the same 20{,}198 to 34{,}133 stamps per task. The agree column
is the fraction of stamps on which the unstable flag matches the
$\sigma_u{=}0.05$ labeling, given as the range over the two off-centre
settings.}
\label{tab:sigmasweep}
\end{table}

\section{Threshold Selection}
\label{app:threshold}

Section~\ref{sec:claim3} states that setting the firing threshold for a learned
head is most of the problem. Table~\ref{tab:thresh} gives the numbers behind
that claim.

\begin{table}[h]
\centering
\small
\setlength{\tabcolsep}{4pt}
\begin{tabular}{lccccc}
\toprule
& & \multicolumn{2}{c}{label-space $\delta$} & \multicolumn{2}{c}{online} \\
\cmidrule(lr){3-4} \cmidrule(lr){5-6}
Task & $\delta$ & $\max \hat\lambda$ & \% over & target & thresh \\
\midrule
OpenCabinet & .085 & .097 & 0.02\% & .06 & .045 \\
OpenDrawer & .076 & .058 & 0.00\% & .14 & .021 \\
PickPlaceToCabinet & .065 & .076 & 0.02\% & .11 & .030 \\
TurnOnSinkFaucet & .058 & .064 & 0.12\% & .12 & .037 \\
\bottomrule
\end{tabular}
\caption{Why the obvious threshold does not work, h1 head. Applying the
label-space $\delta$ to the head's predicted $\hat\lambda$ fires on 0.12
percent of held-out stamps at best and on none at all for OpenDrawer, because
the regression head compresses its outputs toward the mean: the maximum
$\hat\lambda$ it ever emits barely reaches $\delta$. The online columns give
the deployment calibration we adopt, where the threshold is the quantile of a
shadow run's score distribution that reproduces the oracle's firing rate for
that task. Target and achieved rates agree to within 0.1 percentage point on
all eight predictor cells.}
\label{tab:thresh}
\end{table}

Three points are worth separating. First, the failure is not a detection
failure: AUROC is identical under all three threshold choices, because AUROC is
invariant to any monotone rescaling of the score. An offline metric cannot see
this problem at all. Second, calibrating on held-out demonstrations fixes the
scale but not the distribution, and gives online firing rates from 2 to 29
percent against a 4 percent target, because the policy visits states unlike the
demonstrations. Third, the deployment calibration we adopt requires a shadow
run per task and predictor, which is a real cost and one the oracle does not
pay.

\section{Predictor Architecture and Training}
\label{app:predictor}

\paragraph{Features.}
Frozen V-JEPA 2 ViT-L (\texttt{vjepa2-vitl-fpc64-256}), never fine-tuned. For a
stamp at time $t$ we take the 16 camera frames ending at $t$, padded at episode
start, resize to 256, and keep a pooled token grid of 8 temporal positions
$\times$ $4{\times}4$ spatial positions $\times$ 1024 channels, that is 128
tokens per stamp, stored in fp16. Proprioception is end-effector position,
quaternion and gripper joint positions over the same 16-step window.

\paragraph{Head.}
A learned query attends over the token grid with 8-head multi-head attention,
followed by a LayerNorm. Proprioception passes through a single-layer GRU with
64 hidden units. The pooled visual vector and the final GRU state are
concatenated and passed through a two-layer MLP, hidden width 256, GELU,
dropout 0.1, to two scalar outputs: a regression estimate of $\lambda$ and a
classification logit for $\lambda > \delta$. The h4 variant is the same network
reading four clips spaced 16 frames apart with a learned per-clip position
embedding added before attention; h1 is the same code path with one clip.

\paragraph{Training.}
AdamW, learning rate $10^{-3}$, weight decay $10^{-4}$, cosine annealing over
40 epochs, batch size 128. The loss is a Huber loss with $\delta_{\text{H}}{=}0.1$
on $\lambda$ over all stamps plus a positive-weighted binary cross-entropy on
confident stamps only, that is those with $\lvert\lambda\rvert > \delta$, since
the deadband label is ambiguous by construction. The split is by demonstration,
not by stamp, with 20 percent of demonstrations held out, so no frame from a
validation episode is seen in training. Three seeds per configuration; the seed
controls initialization, shuffling and the train/validation split together.

\begin{table}[h]
\centering
\small
\begin{tabular}{lcccc}
\toprule
& \multicolumn{2}{c}{AUROC, confident} & \multicolumn{2}{c}{Spearman with $\lambda$} \\
\cmidrule(lr){2-3} \cmidrule(lr){4-5}
Task & h1 & h4 & h1 & h4 \\
\midrule
OpenCabinet & .885 $\pm$ .021 & .858 $\pm$ .014 & .43 & .34 \\
OpenDrawer & .863 $\pm$ .023 & .838 $\pm$ .010 & .24 & .16 \\
PickPlaceToCabinet & .895 $\pm$ .016 & .887 $\pm$ .006 & .40 & .31 \\
TurnOnSinkFaucet & .902 $\pm$ .022 & .895 $\pm$ .012 & .33 & .28 \\
\bottomrule
\end{tabular}
\caption{Predictor quality on held-out demonstrations, mean and standard
deviation over three seeds. These are the four RoboCasa rows of
Table~\ref{tab:claim2} with the regression quality added, which is the
per-stamp Spearman correlation between predicted and measured $\lambda$ rather
than a ranking of the unstable stratum. The single-clip head is ahead of the
four-clip head on every task and on both metrics, though the AUROC gap exceeds
the seed spread on only two of the four tasks, OpenCabinet and OpenDrawer; the
Spearman column is the clearer separation. Training sets are 16{,}090 to
27{,}974 stamps and validation sets 4{,}108 to 6{,}159 stamps.}
\label{tab:predfull}
\end{table}

\section{Compute}
\label{app:compute}

Labeling is the dominant cost and is CPU-only: no rendering, no network, no
gradient. One RoboCasa task, 200 demonstrations and 20{,}198 to 34{,}133
labeled stamps, takes 2.8 to 6.5 hours on 24 worker processes. Each stamp costs
$(N{+}1) \times K = 216$ simulator steps plus two state restores, so the cost
scales linearly in stamps, in $N$ and in $K$, and the stride is the only knob
that changes stamp count. Feature extraction is one forward pass of a frozen
ViT-L per stamp on one GPU. Head training is minutes on one GPU, the heads
being small relative to the frozen encoder. The simulator oracle of Claim 2 has
the same per-decision cost as one label, 216 steps, which is why it is a bound
rather than a deployable controller.
